\section{The restricted medium convolution estimate}
\label{app:medium}

We prove Proposition~\ref{prop:current-medium}. The convolution lengths
are $U,V$, with $x=UV$, and $\mathfrak{a}$ denotes the nonzero residue.
The even target is denoted by $N$ in the main text.
Put $\vartheta=7/32$ and
\[
 z=x^{1/(\log\log x)^3},\qquad y=x^{1/\log\log x}.
\]
The coefficients $\alpha_n$ are fixed throughout the modulus sums and
satisfy the auxiliary-coprimality Siegel--Walfisz condition
\eqref{eq:current-aux-sw}.
All constants in this condition are uniform over the coefficient family.
We use $e(t)=\exp(2\pi it)$; an inverse in an exponential is always taken
modulo its displayed denominator.

\subsection{A smooth majorant and divisor estimates}

\begin{lemma}\label{ana:majorant}
There is a fixed function $\psi\in C_c^\infty((0,\infty))$ such that
$0\le\psi\le1$, $\psi=1$ on $[1,2]$, $\supp\psi\subset[1/2,5/2]$, and
\[
 \|\psi^{(j)}\|_\infty\le C^j(j!)^2,
 \qquad |\widehat\psi(t)|\ll \exp(-c\sqrt{|t|}).
\]
Consequently, for polynomially bounded positive $\ell,V$ and
$H\ge C_1\ell(\log x)^5/V$,
\begin{equation}
 \frac V\ell\sum_{|h|>H}
       \left|\widehat\psi\left(\frac{hV}{\ell}\right)\right|
       \ll_D x^{-D}
 \label{ana:poisson-tail}
\end{equation}
for every fixed $D>0$.
\end{lemma}

\begin{proof}
Choose $\eta_j=c_0/j^2$ with $\sum \eta_j=1/2$, let $u_j$ be the uniform
probability measure on $[-\eta_j/2,\eta_j/2]$, and put
$\phi=u_1*u_2*\cdots$. Its support is $[-1/4,1/4]$. Differentiating
the first $j$ factors as signed measures gives total variation at most
$\prod_{i\le j}(2/\eta_i)=C^j(j!)^2$. The remaining convolution is a
probability measure; retaining further undifferentiated factors also
gives a continuous density at each fixed derivative order. The function
$\psi=\one_{[3/4,9/4]}*\phi$ has the asserted support and derivative
bounds. Integration by parts $j$ times, with $j$ a sufficiently small
multiple of $\sqrt{|t|}$, gives the Fourier estimate. A decreasing-integral
comparison bounds the tail in \eqref{ana:poisson-tail} by a polynomial
factor times $\exp(-c_1(\log x)^{5/2})$. If $H<1$, every omitted integer
already has $|h|V/\ell\gg(\log x)^5$, and the same argument applies.
\end{proof}

Write $n=\log_xU$, $q=\log_xQ_1$, $u=\log_xQ_2$ and $v=\log_xQ_3$.
The hypotheses of Proposition~\ref{prop:current-medium} imply
\begin{equation}
 n>\sigma,\qquad n\le\frac12-4\sigma,\qquad
 q+u+v<\frac7{10}-\frac{13\sigma}2.
 \label{ana:lengths}
\end{equation}
In particular $Q_1Q_2Q_3\le x^{4/5}$ for large $x$.
Split $q_2=u_0s$ and $q_3=v_0t$ into their $z$-rough and $z$-smooth
parts. The terms $s>y$ or $t>y$ have arbitrarily large logarithmic
saving. To see this, first collapse the modulus coefficient with its
fixed divisor bound. The primitive progression bound for divisor-bounded
functions~\cite{Shiu}, followed by Rankin's method with exponent $1/\log z$, gives
\begin{equation}
 \sum_{\substack{\ell\le x^{4/5}\\ \operatorname{sm}(\ell;z)\ge y}}
       \frac{\tau(\ell)^C}{\ell}
 \ll \exp\bigl(- (\log\log x)^2+O_C(\log\log x)\bigr).
 \label{ana:smooth-tail}
\end{equation}
Here $\operatorname{sm}(\ell;z)$ is the largest divisor of $\ell$ all of
whose prime factors are at most $z$. The bound is used with divisor
means, rather than with a pointwise factor $x^\epsilon$.

We also use the following translated progression estimate. For fixed
$C,K,\delta>0$, an interval $I\subset[-x^K,x^K]$ of length at most $L$
and $f\le Lx^{-\delta}$ satisfy
\begin{equation}
 \sum_{\substack{v\in I\\v\equiv r\pmod f\\v\ne0}}
        \tau(|v|)^C
       \ll \frac Lf\tau(f)^C(\log x)^{O(1)}.
 \label{ana:shifted-divisor}
\end{equation}
Indeed, put $g=(r,f)$ and $v=gw$. The resulting progression modulo
$f/g$ is primitive and its length-to-modulus ratio is $L/f\ge x^\delta$.
The range $|w|\le(f/g)x^{\delta/2}$ is bounded by a fixed divisor
estimate. Split the rest by sign and into dyadic intervals, subdividing
them into lengths comparable to the smaller of their scale and $L/g$.
The primitive divisor bound applies on each interval with fixed positive
length and modulus margins; summation proves \eqref{ana:shifted-divisor}.
Cauchy's inequality gives its mixed version for
$\tau(v)^B\tau(|v-\mathfrak{a}|)^C$, with $v=\mathfrak{a}$ excluded.

On a dyadic block put $d=st\asymp D$, $w=u_0v_0\asymp\mathcal R$ and
define
\begin{equation}
 \Lambda_d(w)=\sum_{\substack{st=d\\u_0v_0=w}}
       \gamma_2(u_0s)\gamma_3(v_0t)
       \one_{\text{fixed individual intervals and supports}}.
 \label{ana:Lambda}
\end{equation}
Thus $\Lambda_d$ is independent of $q_1$,
$|\Lambda_d(w)|\le\tau(dw)$, $D\le y^2=x^{o(1)}$, and
$D\mathcal R\asymp Q_2Q_3$. With $Q=Q_1$, $Y=QD$ and
$\delta=\sigma/3$, \eqref{ana:lengths} implies
\begin{equation}
 U,V\ge x^\delta,\qquad Y\le Ux^{-\delta},\qquad
 Y\mathcal R\le x^{1-\delta}.
 \label{ana:dispersion-lengths}
\end{equation}
We write $\mathcal S$ for the original factorized discrepancy sum
restricted to this block. All values of $d\asymp D$ remain in this
single sum before Cauchy's inequality is applied.

We reduce to one-bounded $\beta$ before applying dispersion. For
divisor-bounded $\beta$, retain $\tau(m)\le(\log x)^C$ and normalize
by $(\log x)^{BC}$. The discarded coefficient at $v=nm$ is at most
$(\log x)^{-C}\tau(v)^{B+2}$. For $v\ne\mathfrak a$, summing the
moduli costs at most $\tau(|v-\mathfrak a|)^3$; fixed shifted divisor
moments therefore give $O(x(\log x)^{E-C})$, with $E$ independent
of $C$. The exact product $v=\mathfrak a$ contributes separately at
most $x^{4/5+o(1)}$. Ordinary divisor means and the reciprocal
totient sum control the mean term. Choose $C$ first, then request
sufficient logarithmic accuracy for the normalized estimate.

\subsection{Dispersion and the principal variance}

The dispersion step separates a common principal term from the
nonzero frequencies. A remaining nonnegative principal variance is
estimated by the auxiliary Siegel--Walfisz hypothesis before the
exponential sums are treated.

Set $\ell=qd$, $\eta_r=\Lambda_d(r)$ and $h_0=\int\psi$. For fixed
$(q,d)$ define
\begin{align*}
 B_r&=\sum_{(n,\ell r)=1}\alpha_n,
 &A_r(b)&=\sum_{\substack{n\equiv b\pmod\ell\\(n,r)=1}}\alpha_n,\\
 C_r(m)&=\frac{\one_{(m,\ell r)=1}}{\varphi(\ell r)}B_r,
 &X_r(m)&=\sum_{(n,\ell r)=1}\alpha_n
       \bigl(\one_{mn\equiv \mathfrak{a}\pmod{\ell r}}-\one_{mn=\mathfrak{a}}\bigr).
\end{align*}
All $n$ sums retain the original interval. Remove $mn=\mathfrak{a}$ before applying
Cauchy's inequality. The first-moment cost is
$Y\mathcal R x^{o(1)}$, which is power-small by
\eqref{ana:dispersion-lengths}. Nonnegative majorants for $q$ and $m$
then give a Cauchy multiplier $VY(\log x)^C$ and the exact variance
\begin{align}
 \mathcal V'&=S_1-2\Re S_2'+S_3',\notag\\
 S_1&=\sum_{q,d}\psi(q/Q)\sum_m\psi(m/V)
                 \left|\sum_r\eta_rC_r(m)\right|^2,\notag\\
 S_2'&=\sum_{q,d}\psi(q/Q)\sum_m\psi(m/V)
       \left(\sum_r\eta_rX_r(m)\right)
       \overline{\left(\sum_s\eta_sC_s(m)\right)},\notag\\
 S_3'&=\sum_{q,d}\psi(q/Q)\sum_m\psi(m/V)
                 \left|\sum_r\eta_rX_r(m)\right|^2.
 \label{ana:variance}
\end{align}

The common principal term is governed by the kernel
\[
 c_\ell(r,s)=\frac{\varphi(\ell rs)}{
       \ell rs\varphi(\ell r)\varphi(\ell s)}.
\]
If $L$ is any common multiple of $\ell$ and the finitely many supported
$r$, then
\begin{equation}
 c_\ell(r,s)=\frac1L\sum_{v\bmod L}
 \frac{\one_{(v,\ell r)=1}}{\varphi(\ell r)}
 \frac{\one_{(v,\ell s)=1}}{\varphi(\ell s)}.
 \label{ana:gram}
\end{equation}
This is a positive Gram kernel. The size of $L$ is irrelevant because it
is used only in the finite identity. Smooth coprime counting yields
\begin{equation}
 S_1=\mathcal M+O\left(\frac{U^2(\log x)^C}{Y}\right),\qquad
 \mathcal M=Vh_0\sum_{q,d}\psi(q/Q)
        \sum_{r,s}\eta_r\overline{\eta_s}
               c_\ell(r,s)B_r\overline{B_s}.
 \label{ana:mean}
\end{equation}

For $S_2'$, remove $(m,s)>1$ while $mn\ne \mathfrak{a}$. A rough common factor
$f>z$ gives factors $\mathcal R/f$ and $x/f$; applying
\eqref{ana:shifted-divisor} and summing
$\sum_{f>z}\tau(f)^C/f^2$ bounds this error by
$VU^2(\log x)^C/(Yz)$. Only then restore $mn=\mathfrak{a}$ for Poisson summation;
the cost is $\mathcal RUx^{o(1)}$. The zero frequency has coefficient
$1/(\ell r\varphi(\ell s))$, whose ratio to $c_\ell(r,s)$ is
$1+O((\log x)/z)$. Consequently
\begin{equation}
 S_2'=\mathcal M+VE_1+
 O_A\left(\frac{VU^2}{Y(\log x)^A}+\mathcal RUx^{o(1)}\right),
 \label{ana:cross}
\end{equation}
where, with the original coprimalities retained,
\begin{align}
 E_1={}&\sum_{q,d}\psi(q/Q)\sum_{r,s}
     \frac{\eta_r\overline{\eta_s}}{\ell r\varphi(\ell s)}
 \sum_{\substack{n_1,n_2\\(n_1,\ell r)=(n_2,\ell s)=1}}
       \alpha_{n_1}\overline{\alpha_{n_2}}\notag\\
 &\hspace{12mm}\times
 \sum_{0<|h|\le H_1}\widehat\psi\left(\frac{hV}{\ell r}\right)
               e\left(\frac{\mathfrak{a}h\overline{n_1}}{\ell r}\right),
 \qquad H_1\asymp\frac{Y\mathcal R}{V}(\log x)^5.
 \label{ana:E1-definition}
\end{align}
Lemma~\ref{ana:majorant} permits this Fourier truncation.

For $S_3'$, retain both exact-product exclusions during all divisor
cutoffs and common-factor removals. The clean set has
\begin{equation}
 \begin{gathered}
 (r,s)=(s,\ell)=1,\quad (n_1,\ell r)=(n_2,\ell s)=1,
 \quad(n_1,n_2)=1,\\
 n_1\equiv n_2\pmod\ell,\qquad
 |n_1-n_2|\ge U(\log x)^{-C_{\rm sep}}.
 \end{gathered}
 \label{ana:clean-set}
\end{equation}
The removed rough common factors cost
$VU^2(\log x)^C/(Yz)+x(\log x)^C$; the close pairs cost
$VU^2(\log x)^{C-C_{\rm sep}}/Y+x(\log x)^C$.
These estimates follow respectively by counting two multiples or a
multiple and a primitive progression; their additional unit terms are
absorbed using $Y/U\le x^{-\delta}$. Choose the divisor truncation and
then $C_{\rm sep}$ sufficiently large.

Now $n_1\ne n_2$. Restoring one equality $mn_i=\mathfrak{a}$ costs only
$\mathcal R(U+Y)x^{o(1)}=\mathcal RUx^{o(1)}$: fixing its
factorization leaves $O(U/Y+1)$ choices for the other $n$, and the other
rough modulus divides a nonzero shifted product. Both equalities cannot
hold. Apply CRT and Poisson only after this restoration. On
\eqref{ana:clean-set},
\[
 \varphi(\ell)c_\ell(r,s)=\frac1{\ell rs}.
\]
Put $d_r(b)=A_r(b)-B_r/\varphi(\ell)$. Since $\sum_b^*A_r(b)=B_r$,
the zero frequency is $\mathcal M+\mathcal P$, where
\begin{equation}
 \mathcal P=Vh_0\sum_{q,d}\psi(q/Q)\varphi(\ell)
   \sum_{r,s}\eta_r\overline{\eta_s}c_\ell(r,s)
                     \sum_b^*d_r(b)\overline{d_s(b)}\ge0.
 \label{ana:principal-variance}
\end{equation}
Positivity follows from \eqref{ana:gram}. In particular this principal
variance must be bounded; it does not cancel identically.

The auxiliary Siegel--Walfisz hypothesis and the primitive character
large sieve give, for every fixed $A,w>0$,
\begin{equation}
 \sum_{\ell\le8Y}\tau(\ell)^w\sum_b^*|d_r(b)|^2
       \ll_A\frac{\tau(r)^{2C_0}U^2}{(\log x)^A}.
 \label{ana:filtered-BDH}
\end{equation}
For clarity, normalize $\alpha_n\one_{(n,r)=1}$ by $\tau(r)^{C_0}$.
Its auxiliary Siegel--Walfisz constants remain uniform because
$\tau(kr)\le\tau(k)\tau(r)$. Character orthogonality expresses the
unweighted variance through nonprincipal characters. Write each modulus
as $fk$, where $f>1$ is the primitive conductor, and keep $(n,k)=1$
in the coefficient before applying the large sieve. On a dyadic
conductor interval $f\asymp R_0$, its bound at fixed $k$ is
$O(U(U/R_0+R_0)/\varphi(k))$. Summing conductors
$f>(\log x)^b$ and inducing factors gives
$O(U^2(\log x)^{2-b}+UY)$. Small conductors are bounded by the
auxiliary Siegel--Walfisz condition with its $\tau(k)^{C_0}$ cost;
their total is $O(U^2(\log x)^{3b+C_1-2B})$.
Choose $b$, and then $B$, to obtain any prescribed logarithmic saving.
The weight $\tau(\ell)^w$ follows by Cauchy's inequality, ordinary
divisor means and the elementary bound for the variance
$O(U^2/\ell)$.

Using \eqref{ana:gram}, Cauchy over the $r$ variables and
$\varphi(\ell r)^{-1}\ll(\log x)/(Y\mathcal R)$ gives, for a fixed
divisor order $B_0$,
\begin{align*}
 \mathcal P\ll{}&\frac{V(\log x)^2}{Y\mathcal R}
  \sum_{r\asymp\mathcal R}\tau(r)^{2B_0}
  \sum_{\ell\le8Y}\tau(\ell)^{2B_0+1}\sum_b^*|d_r(b)|^2\\
 \ll_A{}&\frac{VU^2}{Y(\log x)^A}.
\end{align*}
The factor $\tau(r)^{2C_0}$ in \eqref{ana:filtered-BDH} is summed by a
divisor mean; it is not included in an unspecified uniform constant.
Combining these estimates with \eqref{ana:variance} gives
\begin{equation}
 |\mathcal S|\ll_A\frac{x}{(\log x)^A}
  +V\sqrt Y(\log x)^C\bigl(\sqrt{|E_1|}+\sqrt{|E_2|}\bigr),
 \label{ana:dispersion-bound}
\end{equation}
where $E_2$ is the nonzero CRT frequency sum on
\eqref{ana:clean-set}. The restored single-zero term contributes at
most $\sqrt{xY\mathcal R}\,x^{o(1)}$, which is power-small.

\subsection{The first exponential sum}

For a fixed smooth $w$, Poisson summation and the complete Kloosterman
Weil bound give, for every integer $A$,
\begin{equation}
 \left|\sum_{(n,\ell)=1}w(n/U)e(A\overline n/\ell)\right|
 \ll_w \tau(\ell)^2\sqrt{\ell}+\frac U{\ell}(A,\ell).
 \label{ana:inverse-completion}
\end{equation}
The zero Fourier frequency is a Ramanujan sum and gives the second
term. For the other frequencies use
\[
 \sqrt{(k,A,\ell)}\le\sum_{\substack{d\mid(A,\ell)\\d\mid k}}\sqrt d,
 \qquad
 \sum_{\substack{k\ne0\\d\mid k}}(1+|k|/V)^{-2}\le\frac{2V}{d}.
\]
This proves \eqref{ana:inverse-completion} also when the frequency
length is below one.

Fix $d$ in \eqref{ana:E1-definition}, put $Q'=dQ$ and restrict $q'$
to multiples of $d$. This restriction is permitted inside the bounded
$q'$-dependent coefficient. The following calculation applies with
$Q'$ denoted temporarily by $Q$ and $\mathcal R$ by $R$.
Absorb the bounded smooth factor into the $r,h$ coefficient and apply
Cauchy while retaining the original unrestricted $r,s$ support. Expansion
of the square gives
\begin{equation}
 |E_1(d)|^2\ll x^\epsilon\frac{U^3}{Q^3R^2}\,T,
 \label{ana:E1-cauchy}
\end{equation}
where
\[
 T=\sum_{\substack{q\asymp Q, r,r'\asymp R\\(\mathfrak{a},qrr')=1}}
       \sum_{0<|h|,|h'|\le H_1}
 \left|\sum_{(n,qrr')=1}w(n/U)
 e\left(\frac{\mathfrak{a}(hr'-h'r)\overline n}{qrr'}\right)\right|.
\]
Put $j=hr'-h'r$. The $j=0$ tuples number $RH_1x^\epsilon$, since
$hr'\ne0$ and $h'r=hr'$ is a divisor equation. Their contribution is
$\ll x^\epsilon QRH_1U$. For $j\ne0$, the square-root term in
\eqref{ana:inverse-completion} contributes
$\ll x^\epsilon Q^{3/2}R^3H_1^2$.
For its Ramanujan term write $r=gu$, $r'=gv$, $(u,v)=1$. Then
\[
 (j,rr')\le g^2(h,r)(h',r'),\qquad
 \sum_{q\asymp Q}(j,q)\le3Q\tau(|j|),
\]
and
\[
 \sum_{0<|h|\le H_1}(h,r)\le2H_1\tau(r),\qquad
 \sum_{r,r'\asymp R}(r,r')^2\ll R^3.
\]
Because $(\mathfrak{a},qrr')=1$, $\mathfrak{a}$ disappears from these gcds. The total
Ramanujan contribution to $T$ is $\ll x^\epsilon UH_1^2R$.
As $H_1\ll QRU(\log x)^5/x$, \eqref{ana:E1-cauchy} becomes
\begin{equation}
 |E_1(d)|^2\ll x^\epsilon\frac{U^4}{Q^2}
  \left(\frac{UQ}{x}+\frac{UQ^{5/2}R^3}{x^2}
                         +\frac{U^2QR}{x^2}\right).
 \label{ana:E1-polynomial}
\end{equation}

Return to $Q=Q_1$ and the logarithmic variables in
\eqref{ana:lengths}. The required gaps follow from
\begin{align*}
 n+q&<1-9\sigma,\\
 2n+q+u+v&<\frac{17}{10}-\frac{29\sigma}2,\\
 n+\frac52q+3u+3v
 &=\frac{35}{16}(q-n)+\frac{23}{16}(2n+u+2v)
   +\frac5{16}(n+q+5u+2v)-\frac v2\\
 &<\frac{1021}{512}-\frac{147\sigma}{8}+O(1/\log x).
\end{align*}
The last inequality uses $|\mathfrak{a}|\ge c_0x$ and $\vartheta=7/32$.
The divisor bound for $\Lambda_d$ and $D=x^{o(1)}$ can now be paid
inside a power saving. Sum $U^2/(dQ)$ over $d\asymp D$ using
$\sum_{d\asymp D}1/d\ll1$. With room to absorb one further $D$,
\begin{equation}
 |E_1|\ll\frac{U^2}{DQ}x^{-\sigma/2}.
 \label{ana:E1-final}
\end{equation}

\subsection{The second exponential sum and the reciprocity phase}

Write $E_2=\sum_{d\asymp D}E_3(d)/d$. For fixed $d$, $E_3(d)$ is the
sum over $q$, $w_1,w_2\asymp\mathcal R$, $n_1,n_2\asymp U$ and
$0<|h|\le H$, where $H\asymp QD\mathcal R^2(\log x)^5/V$, of
\begin{align}
 &\frac{\Lambda_d(w_1)\overline{\Lambda_d(w_2)}
               \alpha_{n_1}\overline{\alpha_{n_2}}}{qw_1w_2}
 \psi(q/Q)\widehat\psi\left(\frac{hV}{qdw_1w_2}\right)\notag\\
 &\hspace{8mm}\times
 e\left(\frac{\mathfrak{a}h\overline{n_1w_2}}{qdw_1}
          +\frac{\mathfrak{a}h\overline{n_2qdw_1}}{w_2}\right).
 \label{ana:E3-definition}
\end{align}
The restrictions are
\[
 \begin{gathered}
 (q,\mathfrak{a})=1,\quad(w_1,\mathfrak{a}w_2)=1,\quad(w_2,\mathfrak{a}qdw_1)=1,\\
 n_1\equiv n_2\pmod{qd},\quad
 (n_1,n_2qdw_1)=(n_2,n_1qdw_2)=1,\quad
 |n_1-n_2|\ge U(\log x)^{-C_{\rm sep}}.
 \end{gathered}
\]
There is no independent bounded $q$ coefficient after the outer Cauchy
step. The denominator in \eqref{ana:E3-definition} is $qw_1w_2$;
the factor $1/d$ occurs outside $E_3(d)$.

Exact reciprocity gives
\begin{align}
 &e\left(\frac{\mathfrak{a}h\overline{n_1w_2}}{qdw_1}
          +\frac{\mathfrak{a}h\overline{n_2qdw_1}}{w_2}\right)\notag\\
 &\qquad=
 e\left(\frac{\mathfrak{a}h(n_1-n_2)\overline{n_2qdw_1}}{n_1w_2}\right)
 e\left(\frac{\mathfrak{a}h}{qdw_1w_2n_1}\right).
 \label{eq:current-retained-phase}
\end{align}
All inverses are justified by the displayed coprimalities. Split $h$
by sign and the $q$ support into a bounded number of intervals on scales
$q_0\asymp Q$. Normalize the five variables
$(q,w_1,w_2,n_1,h)$ by $(q_0,\mathcal R,\mathcal R,U,H)$.
The last phase has amplitude
\[
 \frac{\mathfrak{a}}{x}\frac Q{q_0}\frac Dd(\log x)^5.
\]
Mixed differentiation at most once in each coordinate therefore bounds
the total five-dimensional variation by
\begin{equation}
 \bigl(1+(\log x)^5+(|\mathfrak{a}|/x)(\log x)^5\bigr)^5
       =x^{5\kappa+o(1)}.
 \label{ana:phase-variation}
\end{equation}
Repeated partial summation leaves anchored rectangular sums. Their
endpoints are fixed before estimation, and their $h$ coefficient is a
single fixed signed-interval indicator $\omega_h$, $|\omega_h|\le1$.
It has no dependence on any remaining arithmetic variable.

Partition the rough coprime pairs $(n_1w_2,n_2)$ by the cross-coprime
partition of \cite[Lemma 13.2]{MaynardI}; its number of parts is
$\exp(O((\log\log x)^4))$. Expand the two copies of
\eqref{ana:Lambda} independently, including both choices of the
factorization $st=d$. Put $R\asymp v_0$, $S\asymp u_0$.
A total-product interval is retained as
\[
 \one_{A_i<r_is_i\le B_i}
       =\one_{r_is_i\le B_i}-\one_{r_is_i\le A_i}.
\]
Thus it becomes a difference of the permitted product upper cutoffs,
without inserting a coupled condition into a one-variable coefficient.
For a bilinear form, complex polarization replaces two different
coefficient sequences by diagonal sequences with at most four times
the diagonal bound. Apply this successively to the two $R$, two $S$
and two $n$ sequences, at total constant cost at most $4^3$.

We give the resulting unweighted sum explicitly. Fix $d$, a
cross-coprime class $\mathcal C$, endpoints $q_0\le q_1\le2q_0$,
two product upper cutoffs $B_1,B_2$, and one-bounded sequences
$\gamma_r,\lambda_s,\alpha_n,\omega_h$. All these data are fixed
before the remaining summations. Let $\mathscr C$ be the set of
tuples $(q,r_1,r_2,s_1,s_2,n_1,n_2)$ satisfying
\begin{equation}\label{ana:bare-domain}
\begin{gathered}
 q_0\le q\le q_1,\quad (q,\mathfrak{a})=1,\quad
 r_i\asymp R,\quad s_i\asymp S,\quad n_i\asymp U,\\
 (r_1s_1,\mathfrak{a}r_2s_2)=(r_2s_2,\mathfrak{a}qdr_1s_1)=1,\quad
 r_1s_1\le B_1,\quad r_2s_2\le B_2,\\
 n_1\equiv n_2\pmod{qd},\quad
 (n_1,n_2qdr_1s_1)=(n_2,n_1qdr_2s_2)=1,\\
 (n_1r_2s_2,n_2)\in\mathcal C,\qquad
 |n_1-n_2|\ge U(\log x)^{-C_{\rm sep}}.
\end{gathered}
\end{equation}
Define
\begin{equation}\label{ana:bare-definition}
\mathcal W_{\rm bare}=
 \sum_{\mathscr C}
 \frac{\gamma_{r_1}\lambda_{s_1}
           \overline{\gamma_{r_2}\lambda_{s_2}}
           \alpha_{n_1}\overline{\alpha_{n_2}}}
      {q r_1r_2s_1s_2}
 \sum_{0<|h|\le H}\omega_h
 e\left(\frac{\mathfrak{a}h(n_1-n_2)\overline{n_2r_1s_1dq}}
               {n_1r_2s_2}\right).
\end{equation}
There is no smooth reciprocity phase in this definition.
The preceding endpoint separation and polarization give
\begin{equation}\label{ana:E3-bare-reduction}
 |E_3(d)|\ll x^{5\kappa+o(1)}
       \sup|\mathcal W_{\rm bare}|,
\end{equation}
where the supremum is over the choices specified above. The
$x^{o(1)}$ includes the cross-coprime classes and the two independent
factorizations of $d$, but the five-variable phase loss occurs only
once. It remains to bound \eqref{ana:bare-definition}.

\paragraph{The first Cauchy inequality.}
Put $f=(n_1-n_2)/(dq)$, which is a signed nonzero integer. The new
variable satisfies $0<|f|\le2U/(dq_0)$, and
\begin{equation}\label{ana:f-phase}
 e\left(\frac{\mathfrak{a}h(n_1-n_2)\overline{n_2r_1s_1dq}}
                    {n_1r_2s_2}\right)
 =e\left(\frac{\mathfrak{a}hf\overline{n_2r_1s_1}}
                    {n_1r_2s_2}\right).
\end{equation}
Keep $q_f=(n_1-n_2)/(df)$ integral and in $[q_0,q_1]$ at this stage.
The product cutoff in $s_1$ has the exact finite Fourier expansion
\[
 \one_{r_1s_1\le B_1}=
 \int_0^1D_{r_1}(\xi)e(s_1\xi)\,d\xi,
 \qquad
 D_{r_1}(\xi)=\sum_{1\le j\le B_1/r_1}e(-j\xi).
\]
There is a common nonnegative majorant
$L(\xi)\ll\min(1+B_1/R,1+\|\xi\|^{-1})$ with
$\int_0^1L(\xi)\,d\xi\ll\log x$ and
$|D_{r_1}(\xi)|\le L(\xi)$; set
$c_{r_1}(\xi)=D_{r_1}(\xi)/L(\xi)$.
For fixed $\xi$ write
$\lambda_s'=S\lambda_s e(s\xi)/s$.
These coefficients are one-bounded on $s\asymp S$.
If $\mathscr D$ denotes the conditions in \eqref{ana:bare-domain}
which are independent of $s_1$, after putting $q=q_f$, the exact
Fourier integrand is
\begin{align}\label{ana:W2-definition}
 \mathcal W_2(\xi)={}&\frac1S\sum_{\mathscr D}
 \frac{\gamma_{r_1}c_{r_1}(\xi)
              \overline{\gamma_{r_2}\lambda_{s_2}}
              \alpha_{n_1}\overline{\alpha_{n_2}}}
      {q_f r_1r_2s_2}\notag\\
 &\quad\times
 \sum_{\substack{s_1\asymp S\\(s_1,\mathfrak{a}n_1r_2s_2)=1}}
 \lambda_{s_1}'\sum_{0<|h|\le H}\omega_h
 e\left(\frac{\mathfrak{a}hf\overline{n_2r_1s_1}}{n_1r_2s_2}\right).
\end{align}
In particular
$|\mathcal W_{\rm bare}|\ll(\log x)\sup_\xi|\mathcal W_2(\xi)|$.
For each $f$ the number of pairs $n_1,n_2\asymp U$ with
$n_1\equiv n_2\pmod{df}$ is
$O(U^2/(d|f|)+U)$. Summing over $f$ and then $r_1,r_2,s_2$
gives $O(U^2R^2S\log x)$ outer tuples. Since the outside
coefficient in \eqref{ana:W2-definition} has absolute value at most
$O(1/(q_0R^2S^2))$, Cauchy's inequality gives
\begin{equation}\label{ana:first-Cauchy}
 |\mathcal W_2(\xi)|^2
 \ll\frac{U^2\log x}{q_0^2R^2S^3}\,\mathcal W_3,
\end{equation}
where the following nonnegative sum retains the congruence but
discards other outer restrictions:
\begin{align}\label{ana:W3-definition}
 \mathcal W_3={}&
 \sum_{0<|f|\le2U/(dq_0)}
 \sum_{\substack{n_1,n_2\asymp U\\n_1\equiv n_2\pmod{df}}}
 \sum_{r_1,r_2\asymp R}
 \sum_{\substack{s_2\asymp S\\(n_2r_1,n_1r_2s_2)=1}}\notag\\
 &\quad\times
 \left|\sum_{\substack{s_1\asymp S\\(s_1,\mathfrak{a}n_1r_2s_2)=1}}
 \lambda_{s_1}'\sum_{0<|h|\le H}\omega_h
 e\left(\frac{\mathfrak{a}hf\overline{n_2r_1s_1}}{n_1r_2s_2}\right)
 \right|^2.
\end{align}
Every discarded restriction is outside the squared inner sum. The
cross-coprime property of $\mathcal C$ justifies the displayed
coprimality before the original class restriction is removed.

\paragraph{Completion of the positive square.}
We may now remove the congruence in \eqref{ana:W3-definition}
and enlarge $|f|\le2U/q_0$, because all its summands are nonnegative.
Put $b=n_2r_1$ and $c=n_1r_2s_2$. Their multiplicities are at most
$\tau(b)\tau_3(c)$ and can be bounded before summation.
Split $b,c$ and $|f|$ into dyadic intervals. On each such block
$B\ll UR$, $C\ll URS$, $F\ll U/q_0$, and a fixed smooth
nonnegative majorant $g(b,c)=\psi(b/B)\psi(c/C)$ gives
\begin{equation}\label{ana:second-Cauchy}
 \mathcal W_3\ll x^\epsilon\sum_{B,C,F,\pm}\mathcal W_4,
\end{equation}
where the sign of $f$ is absorbed into $\mathfrak{a}$ and
\begin{equation}\label{ana:W4-definition}
 \mathcal W_4=
 \sum_{\substack{b,c\\(b,c)=1}}g(b,c)
 \sum_{f\asymp F}
 \left|\sum_{\substack{s\asymp S\\(s,\mathfrak{a}c)=1}}
 \lambda_s'\sum_{0<|h|\le H}\omega_h
 e\left(\frac{\mathfrak{a}hf\overline{bs}}c\right)\right|^2.
\end{equation}
There are only logarithmically many dyadic blocks. Combining
\eqref{ana:first-Cauchy} and \eqref{ana:second-Cauchy}, and using
$q_0\asymp Q$, gives the outside factor in the form
\begin{equation}\label{ana:bare-Cauchy-factor}
 |\mathcal W_{\rm bare}|
 \ll x^\epsilon\frac{U}{QRS^{3/2}}
       \sup_{B,C,F}\mathcal W_4^{1/2}.
\end{equation}

Expanding \eqref{ana:W4-definition} and interchanging the two
$s$ labels when necessary yields
\begin{equation}\label{ana:W4-expansion}
 \mathcal W_4=
 \sum_{\substack{s_1,s_2\asymp S\\(s_1s_2,\mathfrak{a})=1}}
 \sum_{0<|h_1|,|h_2|\le H}\lambda_{s_2}'\overline{\lambda_{s_1}'}
       \omega_{h_1}\overline{\omega_{h_2}}
 \sum_{f\asymp F}\sum_{\substack{b,c\\(c,bs_1s_2)=1}}
 g(b,c)e\left(\frac{\mathfrak{a}f(h_1s_1-h_2s_2)\overline{bs_1s_2}}c\right).
\end{equation}
The terms with $h_1s_1=h_2s_2$ have total absolute value at most
$x^\epsilon HSBFC$. For the others set
$z=s_1s_2$, $y=f(h_1s_1-h_2s_2)\ne0$, and define the coefficient
\begin{equation}\label{ana:beta-definition}
 \beta(y,z)=
 \sum_{\substack{s_1s_2=z,\ s_i\asymp S\\
       f\asymp F,\ 0<|h_i|\le H\\
       f(h_1s_1-h_2s_2)=y}}
 \lambda_{s_2}'\overline{\lambda_{s_1}'}
       \omega_{h_1}\overline{\omega_{h_2}}.
\end{equation}
We retain $(z,\mathfrak{a})=1$. The variables $b,c$ do not occur in
\eqref{ana:beta-definition}; the restrictions involving them remain
in the kernel. Consequently the nonzero contribution is exactly
\begin{equation}\label{ana:off-diagonal-kernel}
 \sum_{\substack{y\ne0,z\\(\mathfrak{a},z)=1}}\beta(y,z)
 \sum_{\substack{b,c\\(c,bz)=1}}g(b,c)
                         e(\mathfrak{a}y\overline{zb}/c).
\end{equation}
In particular $|\beta(y,z)|\le b(z,y)$ for the positive tuple count
in the following lemma. This domination also holds after any dyadic
restriction on $y,z$. It is independent of the future Poisson
frequency, which has not yet been introduced.

\begin{lemma}\label{ana:multiplicity}
Let $b(z,y)$ count tuples
$s_i\asymp S$, $f\asymp F$, $0<|h_i|\le H$ with
$z=s_1s_2$ and $y=f(h_1s_1-h_2s_2)\ne0$. Then
\begin{equation}
 \sum_z\sum_{u\ne0}b(z,\nu u)^2
 \ll x^\epsilon(1+H/S)\frac{(HS)^2F}{\nu}
 \label{ana:multiplicity-bound}
\end{equation}
for every $\nu\ge1$ on polynomially bounded scales. The left side
vanishes when $\nu\ge8FHS$.
\end{lemma}

\begin{proof}
At fixed $(z,y)$ there are at most $\tau(z)\tau(|y|)$ decompositions
$(s_1,s_2,f,k)$, with $k=h_1s_1-h_2s_2\ne0$.
Cauchy therefore reduces the square to collisions of $h$ pairs with
the same $(s_1,s_2,k)$, followed by a single $f$ count.
Group according to $c\mid\nu$, $\nu/c\mid k$, with the count of $f$
bounded by $2F/c$. Equal $h$ pairs contribute
$\ll x^\epsilon(HS)^2F/\nu$.
For unequal pairs put $g=(s_1,s_2)$, $s_1=gu$, $s_2=gv$, $(u,v)=1$.
Their differences have the form $(h_1-h_1',h_2-h_2')=(jv,ju)$,
with $0<|j|\le2Hg/S$. If $b=(g,\nu/c)$, the condition on
$k/g$ is divisibility by $\nu/(cb)$. Counting nonzero multiples from
zero, without an additional unit term, gives at most
$O(HScb/(g\nu))$ such values. Factoring the remaining nonzero
integer $h_2v$ costs only a divisor factor. The unequal contribution
is consequently
\[
 \ll x^\epsilon\frac{H^3SF}{\nu}
             \sum_{g\le2S}\frac{(g,\nu)}g
 \ll x^{2\epsilon}\frac{H^3SF}{\nu}.
\]
The last estimate follows from
$\sum_{g\le2S}(g,\nu)/g\le\tau(\nu)(1+\log(2S))$.
Renaming $2\epsilon$ proves the result. The support bound follows from
$0<|y|<8FHS$.
\end{proof}

In the present application
$H/S\le x^{o(1)}(Q/U)D_0<1$, where
$D_0=U^2R^2S/x$. Thus the factor $1+H/S$ is bounded.

\subsection{The reciprocal kernel and the spectral estimate}

We estimate the reciprocal kernel with the coefficient independent
of the future Poisson frequency. This separation makes the same
frequency interval available at every spectral level, as required
by the level-average exceptional estimate.

We use $Y$ for the coefficient length, $Z$ for the level length,
and $P$ for the length of a nonzero Poisson-frequency block in this
subsection. The original convolution lengths remain $U,V$.
For $\beta(u,z)$ supported on $u\asymp Y$, $z\asymp Z$ and
$(\mathfrak{a},z)=1$, put
\begin{equation}
 \mathcal K=\sum_{u,z}\beta(u,z)
   \sum_{\substack{c,b\\(c,bz)=1}}
   \psi(c/C)\psi(b/D)e(\mathfrak{a}u\overline{zb}/c).
 \label{ana:kernel}
\end{equation}
Either sign of $u$ is treated separately. Define
$B^2=\sum_{u,z}|\beta(u,z)|^2$ and
$B_\nu^2=\sum_{u,z}|\beta(\nu u,z)|^2$; further coprimality
indicators or unit phases do not increase these norms.

Poisson summation in $b$ is the exact identity
\[
 \sum_{(b,c)=1}\psi(b/D)e(\mathfrak{a}u\overline z\,\overline b/c)
 =\frac Dc\sum_{m\in\mathbb Z}\widehat\psi(mD/c)
                   S(m\overline z,\mathfrak{a}u;c).
\]
The zero frequency contributes at most
$x^\epsilon D\sqrt{YZ}\,B$, since
$|S(0,\mathfrak{a}u;c)|\le(\mathfrak{a}u,c)$ and
$\sum_{c\asymp C}(\mathfrak{a}u,c)\ll C\tau(|\mathfrak{a}u|)$.
Its square is $x^\epsilon D^2YZB^2$.
For nonzero $m$ truncate at the common value $(C/D)x^\epsilon$ and
split into common signed dyadic intervals. Ordinary smooth Fourier
decay, with a fixed sufficiently large derivative order, controls
the discarded tail.

On an actual interval $m\asymp P$, put
\[
 Y_0=\frac{4\pi\sqrt{|\mathfrak{a}|PY}}{C\sqrt Z},\quad
 v=m/P,\quad w=u/Y,\quad \zeta=z/Z,\quad
 t=\frac{4\pi\sqrt{|\mathfrak{a}|mu}}{c\sqrt z\,Y_0}.
\]
Then $c/C=\sqrt{vw/\zeta}/t$. This change of variables removes $\mathfrak{a}$
from the normalized amplitude. Fourier expansion on a fixed compact
box separates the amplitude into a Bessel test function, a common
factor $\one_I(m)e(m\omega)$ and an arbitrary complex $u,z$
coefficient. Choosing the Fourier decay order above the fixed Bessel
derivative order makes the coefficient norm summable. Thus the
$m$ interval $I$ and phase $\omega$ are common at every level $z$.

The cusp conversion and packet operation are given precisely in
Appendix~\ref{app:classical}. They transform the $u$ coefficient to
\begin{equation}
 \beta_\nu(n,z)=\one_{(n,\mathfrak{a})=1}\beta(\nu n,z)e(-n/z),
 \qquad \nu\mid \mathfrak{a}^\infty,
 \label{ana:packet-coefficient}
\end{equation}
at packet cost $(|\mathfrak{a}|\nu)^{\vartheta/2}$, while leaving the common
$m$ functional unchanged. Let
\[
 P_0=\frac{PC^4Z}{|\mathfrak{a}|PY+C^2Z},\qquad T_0=\frac{|\mathfrak{a}|PY}{C^2}.
\]
The regular spectral contribution and the exceptional estimates yield
\begin{align}
 |\mathcal L_P|^2\ll x^\epsilon P_0\bigg\{
 &\sum_{\nu\mid \mathfrak{a}^\infty}(|\mathfrak{a}|\nu)^\vartheta
       (Y/\nu+Z+T_0)(P+Z+T_0)B_\nu^2\notag\\
 &+(C^2Z)^\vartheta(Y+Z)^{1-\vartheta}
                    (P+Z)^{1-2\vartheta}B^2\bigg\}.
 \label{ana:spectral-polynomial}
\end{align}
Here $\mathcal L_P$ is the separated Kuznetsov sum, before restoring
the Poisson factor $D/C$.

We give the spectral summation leading to this polynomial. Throughout
this step, a regular spectral norm is the positive direct-sum norm
\eqref{ana:regular-norm}: its Maass atoms have weight
$1/\cosh(\pi t_j)$, its holomorphic atoms weight $1/(2\pi)$ after
the stated Fourier normalization, and its continuous part measure
$dt/\pi$. Cauchy's inequality is applied within each component of this
same measure, before the three components are added.
At fixed $z,\nu$, write $A(t),B(t)$ for the $n$- and $m$-side
coefficient arrays on $d\mu_z$, with entries
$\sum_n\beta_\nu(n,z)R_{t,1}(n)$ and
$\sum_{m\in I}e(m\omega)R_{t,\infty}(m)$, respectively.
Here $t$ labels an entry of the regular direct sum, including its
component and basis labels, and $R_{t,\mathfrak c}$ denotes the
corresponding Fourier coefficient in~\eqref{ana:regular-norm}.
Let $A_0,B_0\ge0$ denote the norm-bound amplitudes, and let
$u_0,v_0\ge0$. If the
two cumulative regular squared norms are bounded by
$(T^2+u_0)A_0^2$ and $(T^2+v_0)B_0^2$, the Bessel bounds
\[
 |\mathcal H(t)|\ll\frac{1+|\log Y_0|}{1+Y_0},\qquad
 |\mathcal H(t)|\ll |t|^{-5/2}+Y_0|t|^{-3}
 \quad(|t|\ge\max(1,2Y_0))
\]
give
\[
 \int|\mathcal H(t)A(t)B(t)|\,d\mu_z(t)
 \ll \frac{1+|\log Y_0|}{1+Y_0}
 \sqrt{(1+Y_0^2+u_0)(1+Y_0^2+v_0)}A_0B_0.
\]
Indeed, above $\max(1,2Y_0)$ the cumulative norm product grows on
the $j$th shell by at most $2^{2j}$, leaving the summable factors
$2^{-j/2}+2^{-j}$. The level Cauchy factor is $\sqrt Z$.
With $L=Y/\nu$, this gives the first term of
\eqref{ana:spectral-polynomial}, including the packet cost.

For exceptional forms put $\tau_j=2\sqrt{1/4-\lambda_j}\le\vartheta$,
$X_0=1+Z/L$, and $X=1/Y_0$. Use
\[
 1+X^{\tau_j}\le
 2(1+X/\sqrt{X_0})^{\tau_j}(1+X_0)^{\tau_j/2}.
\]
The fixed-level $n$ norm is
$\ll x^\epsilon(1+L/Z)B_\nu^2$. The common-interval $m$ norm is
\[
 \ll x^\epsilon P(P+Z)
 \left[1+\left\{
  \frac{P(1+X/\sqrt{X_0})^2}{(P+Z)^2}\right\}^{\vartheta}\right].
\]
Multiplying these bounds and the squared packet cost gives a term
absorbed in the regular expression and the second term of
\eqref{ana:spectral-polynomial}. The factor $|\mathfrak{a}|\nu$ cancels only
in this latter term after the two norms are multiplied.
The same argument applies to the holomorphic component with height
$k-1$ and to the continuous component with real height $t$ in
\eqref{ana:regular-norm}. Their measures remain unchanged throughout
the shell summation. On the exceptional spectrum,
$1/\cosh(\pi t_j)=1/\cos(\pi|\operatorname{Im}t_j|)\asymp1$,
since $|\operatorname{Im}t_j|\le7/64$; hence the unweighted
exceptional estimates used above differ only by an absolute constant.

Restoring $D/C$, the elementary inequality
\[
 \frac{Z(L+Z+T_0)(P+Z+T_0)}{Z+T_0}
       \le(L+Z)(P+Z)+ZT_0
\]
and $P\le C/D$ give
\[
 |\mathcal K|^2\ll x^\epsilon\mathcal J^2,
\]
where the complete nonnegative polynomial is
\begin{align}
 \mathcal J^2={}&\sum_{\nu\mid \mathfrak{a}^\infty}(|\mathfrak{a}|\nu)^\vartheta
 \bigl\{C(Y/\nu+Z)(C+DZ)+|\mathfrak{a}|YZ\bigr\}B_\nu^2
 \notag\\
 &+\bigl\{
 C(CD\sqrt Z)^{2\vartheta}(Y+Z)^{1-\vartheta}
                         (C+DZ)^{1-2\vartheta}+D^2YZ\bigr\}B^2.
 \label{ana:returned-polynomial}
\end{align}
The actual cutoff $P\le x^\epsilon C/D$ is accommodated by reducing
the initially chosen $\epsilon$. If $C/D<1/2$, estimates are applied
only on nonempty actual intervals; the real comparison value $C/D$
is not used as the length of a nonexistent integer interval.

\paragraph{Return to the convolution scales.}
Apply \eqref{ana:returned-polynomial} first at the actual dyadic
lengths in \eqref{ana:off-diagonal-kernel}. All its coefficients are
nonnegative and increasing in the four lengths, since
$0\le\vartheta\le1/2$. We may then replace these lengths by
\begin{equation}\label{ana:scale-substitution}
 D=B_*=UR,\quad C=C_*=URS,\quad Z=S^2,\quad
 Y=\frac{UR^2S^3}{V}.
\end{equation}
The accompanying counting parameters satisfy
$F\le U/Q$ and $H\le x^\epsilon H_0$, where
$H_0=QR^2S^2/V$. These selectable powers are included in the final
$x^\epsilon$ after choosing the input exponents smaller. Put
\[
 D_0=\frac{U^2R^2S}{x}=\frac YZ<1.
\]
Then $B_*Z\ge C_*$, $B_*^2<x$, and $C_*B_*Z=xY$.
Lemma~\ref{ana:multiplicity} and
$H/S\le x^\epsilon(Q/U)D_0<1$ give
\begin{equation}\label{ana:norm-insertion}
 B_\nu^2\le x^\epsilon\frac{(HS)^2F}{\nu},\qquad
 B^2\le x^\epsilon(HS)^2F.
\end{equation}
The auxiliary divisor sum satisfies
\[
 \sum_{\nu\mid\mathfrak a^\infty}\nu^{\vartheta-1}
 =\prod_{p\mid\mathfrak a}(1-p^{-(1-\vartheta)})^{-1}
 =x^{o(1)}.
\]
Its logarithm is $O((\log x)^\vartheta)$, by splitting the prime
divisors at $\log x$; this is uniform for polynomially bounded
$\mathfrak a$. Consequently \eqref{ana:returned-polynomial} gives
\begin{equation}\label{ana:J-reduced}
 \mathcal J^2\ll x^\epsilon(HS)^2F
 \left\{
 |\mathfrak a|^\vartheta(C_*B_*Z+|\mathfrak a|Y)Z
 +C_*^{1+2\vartheta}B_*Z^{2-2\vartheta}
 +B_*^2YZ\right\}.
\end{equation}
For clarity, the four monomials on the right yield the following
ratios after multiplication by
\[
 \frac{(H_0S)^2(U/Q)}{U^4R^4S^6}
                  =\frac{Q}{U^3V^2}:
\]
\begin{equation}\label{ana:ratio-table}
\begin{array}{c|c}
 \text{monomial}&\text{normalized contribution}\\ \hline
 |\mathfrak a|^\vartheta C_*B_*Z^2
     &D_1=|\mathfrak a|^\vartheta UQR^2S^5/x^2\\
 |\mathfrak a|^{1+\vartheta}YZ
     &(|\mathfrak a|/x)D_1\\
 C_*^{1+2\vartheta}B_*Z^{2-2\vartheta}
     &D_2=U^{1+2\vartheta}QR^{2+2\vartheta}S^{5-2\vartheta}/x^2\\
 B_*^2YZ&QR^4S^5/V^3\le D_1.
\end{array}
\end{equation}
The last comparison follows from $B_*^2<x$ and $|\mathfrak a|\ge1$.
The zero-difference contribution in \eqref{ana:W4-expansion}
satisfies
\[
 \left(\frac{U}{QRS^{3/2}}\right)^2H_0SB_*\frac UQ C_*
       =\left(\frac{U^2}{Q}\right)^2D_0.
\]
For the nonzero part, \eqref{ana:off-diagonal-kernel} gives
$|\mathcal W_{4,\ne0}|\ll x^\epsilon\mathcal J$.
Combining this square root with that in
\eqref{ana:bare-Cauchy-factor} gives explicitly
\[
 \left(\frac{U}{QRS^{3/2}}\right)^4
 \frac{\mathcal J^2}{(U^2/Q)^4}
 =\frac{\mathcal J^2}{U^4R^4S^6}
 \ll x^\epsilon\bigl\{(1+|\mathfrak a|/x)D_1+D_2\bigr\}.
\]
The two contributions prove
\begin{align}
 |\mathcal W_{\rm bare}|&\ll\frac{U^2}{Q}x^\epsilon
   \left\{D_0^{1/2}
    +\bigl[(1+|\mathfrak{a}|/x)D_1+D_2\bigr]^{1/4}\right\},
 \label{eq:current-bare-bound}\\
 D_0&=\frac{U^2R^2S}{x},\qquad
 D_1=\frac{|\mathfrak{a}|^\vartheta UQR^2S^5}{x^2},\qquad
 D_2=\frac{U^{1+2\vartheta}QR^{2+2\vartheta}S^{5-2\vartheta}}{x^2}.
 \notag
\end{align}

\subsection{Completion of the convolution estimate}

The scale hypotheses give $D_0\ll x^{-8\sigma}$ and
$D_1,D_2\ll x^{-15\sigma}$ on every small-factor block.
Together with $|\mathfrak{a}|/x\le x^\kappa$ and the single variation cost
\eqref{ana:phase-variation}, \eqref{eq:current-bare-bound} gives
\[
 |E_3(d)|\ll\frac{U^2}{Q}
 \left(x^{-4\sigma+5\kappa+\epsilon}
       +x^{-15\sigma/4+21\kappa/4+\epsilon}\right).
\]
All selectable epsilon losses, the vanishing subpowers from the
partitions and one additional normalizing factor $D$ fit inside
$x^{\sigma/10}$. Since $\kappa\le\sigma/100$, both remaining
savings exceed $\sigma/2$. Harmonic summation in $d$ therefore gives
\begin{equation}
 |E_2|\ll\frac{U^2}{DQ}x^{-\sigma/2}.
 \label{ana:E2-final}
\end{equation}
Insert \eqref{ana:E1-final} and \eqref{ana:E2-final} into
\eqref{ana:dispersion-bound}. The exponential-sum contribution is
$O(x^{1-\sigma/4}(\log x)^C)$. The earlier single-zero restoration
is $O(x^{1-\sigma/6+o(1)})$, still power-small.
Choose the auxiliary Siegel--Walfisz and logarithmic precisions to
pay for the finitely many dyadic choices, the fixed logarithmic number
of blocks and the normalization of $\beta$. Equation
\eqref{ana:smooth-tail} deals with the omitted smooth parts.
This proves \eqref{eq:current-medium-conclusion}.

\section{Normalization of the spectral inputs}
\label{app:classical}

We record the normalizations to identify the positive measures used
in the spectral Cauchy inequalities. The regular spectrum and the
exceptional spectrum are treated in their respective ranges.

These inputs are used in \eqref{ana:spectral-polynomial}; the packet coefficient and the common frequency interval are specified in \eqref{ana:packet-coefficient} and the discussion preceding that estimate.

We specify the classical statements used in the preceding proof.
They are applied at trivial character, level $z$, cusps $\infty$ and
$1$, with $(\mathfrak{a},z)=1$. Write
$R(n)=\sqrt{|n|}\rho(n)$ for the Deshouillers--Iwaniec normalization
of the Fourier coefficients. The cusp scaling matrices are
\[
 \sigma_\infty=I,\qquad
 \sigma_1=\begin{pmatrix}\sqrt z&0\\\sqrt z&1/\sqrt z\end{pmatrix}.
\]
The cusp conversion in \cite[Lemma 4.3]{Drappeau} introduces
$e(\mathfrak{a}y/z)$; hence the ordinary Kloosterman coefficient carries
$e(-\mathfrak{a}y/z)$.

For $y=\nu n$, $\nu\mid \mathfrak{a}^\infty$, $(n,\mathfrak{a})=1$, the coprime packet
identities in \cite[Lemmas 3.1 and 3.2]{ABL} give, on every normalized
oldform dilation,
\[
 R_v(\mathfrak{a}\nu n)=\lambda_v(\mathfrak{a}\nu)
                   e\bigl(n(\mathfrak{a}\nu-1)/z\bigr)R_v(n).
\]
Multiplication by $e(-\mathfrak{a}\nu n/z)$ leaves exactly $e(-n/z)$.
The packet coefficients in \cite[(3.3)]{ABL} have Frobenius norm
$z^\epsilon$ after divisor losses. Cauchy's inequality in that packet
and the Kim--Sarnak exponent bound~\cite{KimSarnak} therefore cost
$(|\mathfrak{a}|\nu)^{7/64+\epsilon}$ and give the common coefficient
\eqref{ana:packet-coefficient}. No eigenvalue is assumed common to
all dilations. The holomorphic and continuous quadratic forms in
\cite[Lemma 3.2]{ABL} use the corresponding identities with
$\vartheta=0$. The continuous term is integrated with the formal
Plancherel measure, rather than treated as the norm of an individual
square-integrable Eisenstein series; see also
\cite[Proposition 8.2]{Young}.

For completeness we specify the positive norm used for every regular
spectral Cauchy inequality. Let $\mathfrak c$ be either cusp. In the
notation of \cite[Proposition E]{Pascadi}, write
\begin{align*}
 R^M_{j,\mathfrak c}(n)&=
       \rho^{\rm DI}_{j,\mathfrak c}(n)
       =\sqrt{|n|}\rho^{\rm ABL}_{j,\mathfrak c}(n),\\
 R^H_{jk,\mathfrak c}(n)&=
       \left(\frac{(k-1)!}{(4\pi)^{k-1}}\right)^{1/2}
       n^{-(k-1)/2}\psi_{jk,\mathfrak c}(n),\\
 R^E_{\mathfrak d,\mathfrak c}(n,t)&=
       n^{it}\varphi_{\mathfrak d\mathfrak c}
                       (n;\tfrac12+it).
\end{align*}
The symbol $R_j$ in the packet formulas denotes $R_j^M$.
For $T\ge1$ and $b$ supported on $n\asymp L$, $L\ge1/2$, define
\begin{align}\label{ana:regular-norm}
 \mathfrak Q_{z,\mathfrak c}(T;b)
 ={}&\sum_{\substack{t_j\in\mathbb R\\|t_j|\le T}}
 \frac1{\cosh(\pi t_j)}
 \left|\sum_{n\asymp L}b_nR^M_{j,\mathfrak c}(n)\right|^2\notag\\
 &+\frac1{2\pi}\sum_{\substack{k\in2\mathbb Z_{>0}\\k-1\le T}}
      \sum_j\left|\sum_{n\asymp L}b_nR^H_{jk,\mathfrak c}(n)\right|^2
       \notag\\
 &+\frac1\pi\sum_{\mathfrak d}\int_{-T}^T
      \left|\sum_{n\asymp L}b_n
                 R^E_{\mathfrak d,\mathfrak c}(n,t)\right|^2dt.
\end{align}
This is the direct-sum measure $d\mu_z$ used above. The constants
$1/(2\pi)$ and $1/\pi$ are the constants in the Kuznetsov formula;
the factor $1/\cosh(\pi t_j)$ is part of the Maass measure and is
not absorbed by the change from $\rho^{\rm ABL}$ to $R^M$.
By \cite[Proposition 4.7]{Drappeau} and
\cite[Lemma H]{Pascadi},
\begin{equation}\label{ana:regular-large-sieve}
 \mathfrak Q_{z,\mathfrak c}(T;b)
 \ll (zLT)^\epsilon\left(T^2+\frac Lz\right)\sum_n|b_n|^2.
\end{equation}
The harmless change from $k\le T$ to $k-1\le T$ costs only a fixed
factor. The continuous term in this convention has measure $dt/\pi$,
as in \cite[(3.27)]{Pascadi}; it has no additional hyperbolic factor.
The test-function bounds used with this norm are
\cite[Lemma D]{Pascadi}. For opposite-sign Kuznetsov sums the
holomorphic term is absent, and the same positive norm is still an
upper bound.

Put $\tau_j=2\sqrt{1/4-\lambda_j}\le7/32$ on the exceptional
spectrum. The fixed-level estimate \cite[Theorem A]{Pascadi}, at
cusp $1$, supplies
\begin{equation}
 \sum_{\rm exc}X^{\tau_j}
       \left|\sum_{n\asymp L}b_nR_{j,1}(n)\right|^2
 \ll(zL)^\epsilon(1+L/z)\sum_n|b_n|^2,
 \qquad X\ll\max(1,z/L).
 \label{ana:exceptional-fixed}
\end{equation}
The coefficients may depend on $z$ but are independent of spectral
height. The level-average estimate \cite[Theorem J]{Pascadi} gives
\begin{equation}
 \sum_{z\asymp Z}\sum_{\rm exc}X^{\tau_j}
       \left|\sum_{m\in I}e(m\omega)R_{j,\infty}(m)\right|^2
 \ll(ZV)^\epsilon(Z+V)V,
 \quad X\ll\max(V,Z^2/V),
 \label{ana:exceptional-average}
\end{equation}
where $I\subset(V,2V]$ and $\omega$ are common to every level.
For a partial interval $I$, choose an integer $P>4V$ and expand
$\one_I(m)=\sum_{k=0}^{P-1}c_ke(km/P)$ on residues modulo $P$.
The coefficient norm is $\sum|c_k|\ll\log(2+V)$.
Minkowski in the Hilbert norm indexed by $(z,j)$ and
Theorem J with phase $\omega+k/P$ prove
\eqref{ana:exceptional-average}, with only a logarithmic squared
loss absorbed in $(ZV)^\epsilon$. This argument requires the same
interval at every level. For larger $X$, the bound
$\tau_j\le\vartheta$ gives the continuation factor
$1+(VX/(V+Z)^2)^\vartheta$.

The transformed estimates apply on each nonempty interval of length
at least $1/2$, including an endpoint interval containing $n=1$.
The $m$ interval is common to all levels, and the $n$ coefficient
is independent of spectral height. The condition
$(\mathfrak a,z)=1$ is retained in every packet operation.
These are precisely the interval and coefficient restrictions used
in the preceding spectral estimate.